% =============================================================================
% Section 4: Black Holes in The Interstellar Medium
% Merged from agents 12--17
% Subsections 4.1 through 4.9
% =============================================================================

\section{Black Holes in the Interstellar Medium}
\label{sec:4}

%%% ====================================================================
%%% 4.1  Accretion of Noninteracting Particles onto a Nonmoving Black Hole
%%% ====================================================================

\subsection{Accretion of Noninteracting Particles onto a Nonmoving Black Hole}
\label{sec:4.1}

Consider a black hole at rest in the interstellar medium; and temporarily treat
the interstellar gas as though it were made up of noninteracting particles
[collisions neglected; mean free paths large compared to the region over which
the hole's gravity makes itself felt; $l_\text{fp} \gg 2GM/c^2$; see below]. In the case of a
Schwarzschild hole, all particles with angular momentum per unit mass
$\ell < 2r_g c$  eventually get captured by the hole (see Box~25.6 of MTW or Figure~12
of ZN). Here $r_g = 2GM/c^2$ is the gravitational radius of hole and $M$ is the hole's
mass. If the particle speeds far from the hole are $v_\infty$, then this condition for
capture corresponds to an impact parameter
%
\begin{equation}
b = \ell / v_\infty = 2r_g c/v_\infty\,.
\tag{4.1.1}
\end{equation}
%
Consequently, the ``capture cross section'' of the hole is
%
\begin{equation}
\sigma_\text{capture} = \pi b^2 = 4\pi r_g^2 (c/v_\infty)^2\,.
\tag{4.1.2}
\end{equation}
%
For a Kerr hole the capture cross section is of this same order of magnitude. The
rest mass per unit time crossing inward through a sphere of $r \gg r_\text{capture}$, with
particles directed into the capture region (solid angle $\Delta\Omega = \sigma/r^2$), is
%
\begin{equation}
\dot{M}_0 = \left(\frac{\rho_\text{cap.}}{4\pi}\right) 4\pi r^2 \Delta\Omega = 4\pi r_g^2 \rho_\infty c^2 / v_\infty\,.
\tag{4.1.3}
\end{equation}
%
This is the rate at which the hole accretes rest mass. Rewritten in typical astro\-nomical units, this accretion rate is
%
\begin{equation}
\dot{M}_0 = 10^{11} \left(\frac{\rho_\infty}{10^{-24}~\text{g~cm}^{-3}}\right) \left(\frac{M}{M_\odot}\right)^2 \left(\frac{v_\infty}{10~\text{km~sec}^{-1}}\right)^{-1} \text{``H~II''}
\tag{4.1.3$'$}
\end{equation}
%
\begin{equation}
\frac{d(M_0/M_\odot)}{d(t/10^{10}~\text{yrs})} = 10^{-13} \left(\frac{\rho_\infty}{10^{-24}~\text{g~cm}^{-3}}\right) \left(\frac{M}{M_\odot}\right)^2 \left(\frac{v_\infty}{10~\text{km~sec}^{-1}}\right)^{-1}\,.
\notag
\end{equation}

[The typical densities and speeds of interstellar gas particles in ionized, ``H~II''
regions are $10^{-24}$~g~cm$^{-3}$ and 10~km~sec$^{-1}$; see e.g.\ Kaplan and Pikel'ner (1970)].
Even if 100 per cent of the inflowing rest mass were somehow converted into
outgoing radiation, the total luminosity would be only
%
\begin{equation}
L_\text{max} = \left(10^{24}~\frac{\text{erg}}{\text{sec}}\right) \left(\frac{\rho_\infty}{10^{-24}~\text{g~cm}^{-3}}\right) \left(\frac{M}{M_\odot}\right)^2 \left(\frac{v_\infty}{10~\text{km~sec}^{-1}}\right)^{-1}
\tag{4.1.4}
\end{equation}
%
---a value much too small to be astronomically interesting. Therefore it is fortunate
that the electrons, ions, and magnetic fields of the interstellar gas interact strongly
enough to make the accretion process obey fluid-dynamic laws rather than the
laws of noninteracting particles (see \S4.6, below).

%%% ====================================================================
%%% 4.2  Adiabatic, Hydrodynamic Accretion onto a Nonmoving Black Hole
%%% ====================================================================

\subsection{Adiabatic, Hydrodynamic Accretion onto a Nonmoving Black Hole}
\label{sec:4.2}

Switch, then, from a noninteracting description of accretion to a hydrodynamic
description. Assume, as above, that the black hole is at rest with respect to the gas
and a Schwarzschild hole, so that the accretion is spherically symmetric. In the
hydrodynamic case the accretion rate $\dot{M}_0$ and all other basic characteristics of the
flow are governed by the gravitational field at distances much greater than the
gravitational radius. (This is because the flow at small radii is supersonic and
therefore cannot influence conditions at large radii.) Thus, one can calculate
the mass flow and other quantities at large radii accurately using the Newtonian
theory of gravitation. Phenomena close to the gravitational radius will be treated
later. The motion of the gas will be assumed adiabatic. Deviations from adiabatic
flow due to radiative losses and radiative transport between gas elements can be
taken into account later by suitably modifying the adiabatic index $\Gamma$.
The form of the flow is governed by two fundamental equations: conservation
of rest mass (``continuity equation''), which we write in the form\footnote{Material for this section is drawn largely from Chapter~13 of ZN, and from Shvartsman and Novikov (1971).}
%
\begin{equation}
4\pi r^2 \rho u = \dot{M}_0 = \text{constant, independent of}~r,
\tag{4.2.1}
\end{equation}
%
and the Euler equation
%
\begin{equation}
u\frac{du}{dr} = -\frac{1}{\rho}\frac{dp}{dr} - \frac{GM}{r^2}\,.
\tag{4.2.2}
\end{equation}
%
[Cf.\ eqs.\ (2.5.23) and (2.5.59).]

Here $\dot{M}_0$ is the total rate of accretion of rest mass, $u$ is the radial
velocity, and $M$ is the mass of the hole.  We assume the gas has constant
adiabatic index~$\Gamma$, so that during the accretion the pressure and
density are related by the adiabatic law $p = K\rho^\Gamma$, and the speed of
sound is given by $a = (\Gamma p/\rho)^{1/2}$.  Let the constants $K$ and $\Gamma$
be given.  Our task is to determine the accretion rate~$\dot{M}_0$ and to
compute as well the distributions of density $\rho(r)$ and velocity $u(r)$
in the flow.\footnote{Notice that our notation differs from that in \S2.
Throughout \S4 we denote (for ease of eyesight)

(speed of sound) $= a$, not $c_s$

(density of rest mass) $= \rho$, not $\rho_0$

No confusion is likely, since nowhere in \S4 shall we deal with
4-accelerations or with total mass-energy; and
$\rho = \rho_0(1+\Pi)$, and because $\Pi \ll 1$ in all regions of the accreting gas.
We shall also retain factors of $G$, $c$, and $k$ (i.e.\ use cgs units)
throughout \S4.}

In place of the Euler equation~(4.2.2) we shall use the Bernoulli
equation (2.5.38) rewritten in the form
%
\begin{equation}
\tfrac{1}{2}\,u^2 + \frac{1}{\Gamma - 1}\,a^2 - \frac{GM}{r}
= \text{constant}
= \frac{1}{\Gamma - 1}\,a_\infty^2\,.
\tag{4.2.3}
\end{equation}
%
[Here the enthalpy has been written in the form
$w = \int \rho^{-1}\,dp = a^2/(\Gamma-1)$;
cf.\ eq.\ (2.5.42).]  The constant has been determined by conditions at
infinity, where the gas is at rest.

Rewrite the law of mass conservation~(4.2.1) with density expressed in terms
of sound speed
%
\begin{equation}
u = \frac{\dot{M}_0}{4\pi\rho_\infty\,r^2}
    \left(\frac{a_\infty}{a}\right)^{2/(\Gamma-1)}\!.
\tag{4.2.4}
\end{equation}

The only unknown parameter in the coupled equations~(4.2.3) and~(4.2.4) is
the accretion rate~$\dot{M}_0$.  Once $\dot{M}_0$ is known, one can readily
solve for the radial distributions of velocity, sound speed, and density.

The mass flux is determined in the following way.  Consider the system of
equations~(4.2.3) and~(4.2.4).  In the $u$--$a$ plane, the Bernoulli
equation~(4.2.3) for fixed radius~$r$ defines an ellipse; each value of $r$
corresponds to a different ellipse.  Similarly, the equation of mass
conservation~(4.2.4) for each value of $r$ defines a hyperbola of fractional
power, $u\,a^{2/(\Gamma-1)} = \text{const}$.  Now, let the accretion
rate~$\dot{M}_0$ be chosen arbitrarily.  For every value of $r$,
equations~(4.2.3) and~(4.2.4) are 2 equations with 2 unknowns, $u$ and $a$.
Solving these equations---i.e.\ finding the intersection point of the ellipse
with the hyperbola---gives $u$ and $a$ for that particular radius.  In other
words, in the $u$, $a$ plane the curve~$u(a)$ is determined parametrically by
the intersections of corresponding ellipses and hyperbolae.  It is clear (see
Figure~\ref{fig:4.2.1}) that for every ellipse-hyperbola pair, there are
either two intersection points, or one point of tangency, or no intersection
at all.  If, for a particular chosen~$\dot{M}_0$, there exists any $r$ at
which the two equations are incompatible for the~$\dot{M}_0$, such flow
cannot occur.  Rejecting this case, we have 2 remaining possibilities, for a
given~$\dot{M}_0$ (see Figure~\ref{fig:4.2.2}):

\begin{enumerate}
\item The ellipse and hyperbola intersect twice for every value of~$r$
  (Figure~\ref{fig:4.2.2}a).\\
  In this case we have 2 separate curves, corresponding to 2 families of
  intersection points.

\item The ellipse and hyperbola meet tangentially
  (Figure~\ref{fig:4.2.2}b) at one value of $r$ (call it~$r_s$), at which
  they cross each other; the two families of intersection points $u(a)$
  cross each other at $r = r_s$.  We shall show below that this
  crossing point necessarily lies on the ``bisectrix'' of the graph,
  $u = a$.
\end{enumerate}

In both the cases, 1 and 2, the curves $u(a)$ describe possible flows of gas
in the gravitational field.  The curves begin on the innermost ellipse,
$r = \infty$.  On the upper curve~$u(a)$, at this ellipse we have $a = 0$,
but $u \neq 0$.  These boundary conditions are not of interest for the
accretion problem,\footnote{This curve describes the ``stellar wind'' by
which some stars eject matter into interstellar space; see Chapter~13 of
ZN.} because accretion requires $u_\infty = 0$,
$a_\infty \neq 0$.  The lower curve~$u(a)$ begins at the point
$u_\infty = 0$, $a_\infty \neq 0$---

\begin{figure}[htbp]
\centering
% \includegraphics[width=\columnwidth]{fig_4_2_1}
\fbox{\parbox{0.9\columnwidth}{\centering [Figure 4.2.1 placeholder]\\[6pt]
The $u$, $a$ ``solution plane'' for solving the coupled Bernoulli
equation~(4.2.3) and law of mass conservation~(4.2.4).
Solid curves: ellipses from the Bernoulli equation at several fixed values
of~$r$ ($r = \infty$ is the innermost ellipse).
Dashed curves: hyperbolae from the continuity equation.
Heavy dots mark intersection points; the locus of these intersections
traces the flow solution $u(a)$.
The diagonal $u = a$ is the ``bisectrix''.
Regions above the bisectrix are supersonic; regions below are subsonic.}}
\caption{The $u$, $a$ ``solution plane'' for solving the coupled Bernoulli
equation~(4.2.3) and law of mass conservation~(4.2.4).}
\label{fig:4.2.1}
\end{figure}

to our desired boundary condition.  Hence we shall restrict attention to the
lower curve.

We must still decide which type of lower curve is reasonable---one that
remains always below the bisectrix (case~1), or one that crosses the
bisectrix (case~2).  In case~1 the flow is subsonic at all radii, $u < a$.
But this requires a large back-pressure at all radii to retard the
inflow---back pressure that can never be provided near the gravitational
radius of a black hole.  The gas must cross the gravitational

\begin{figure}[htbp]
\centering
% \includegraphics[width=\columnwidth]{fig_4_2_2}
\fbox{\parbox{0.9\columnwidth}{\centering [Figure 4.2.2 placeholder]\\[6pt]
Possible forms of adiabatic gas flow in a spherical, Newtonian gravitational
field.  Curves~(a) [case~1 in text] correspond to flow that remains subsonic
or supersonic everywhere.  Curves~(b) [case~2 in text] correspond to flow
for which there is a sonic point, $u = a$.}}
\caption{Possible forms of adiabatic gas flow in a spherical, Newtonian
gravitational field.  Curves~(a) [case~1 in text] correspond to flow that
remains subsonic or supersonic everywhere.  Curves~(b) [case~2 in text]
correspond to flow for which there is a sonic point, $u = a$.}
\label{fig:4.2.2}
\end{figure}

\noindent
radius with the speed of light, as measured in the \emph{proper} reference
frame of an observer there; hence, the flow is surely supersonic near the
gravitational radius.

Thus, the only acceptable solution $u(a)$ is the lower curve of case~2
(Figure~\ref{fig:4.2.2}b).  This curve begins at $u = u_\infty = 0$,
$a = a_\infty \neq 0$, and it crosses from the subsonic region into the
supersonic region at $r = r_s$.  This transition to supersonic flow is
crucial to the accretion process.  It governs the accretion rate~$\dot{M}_0$,
in the same manner as the transition to supersonic flow in the throat of a
rocket nozzle.  In both cases one and only one mass flow rate is compatible
with the required transition to supersonic flow.

Let us calculate the required accretion rate~$\dot{M}_0$.  We begin by
calculating the flow velocity at the transition radius (``sonic
radius'')~$r_s$.  To do this, we rewrite the Euler equation~(4.2.2) in the
form
%
\begin{equation}
\frac{du}{dr} = -\,\frac{a^2\,dp}{p\,dr} - \frac{GM}{r^2}\,.
\tag{4.2.5}
\end{equation}

We differentiate the law of mass conservation~(4.2.1) with respect to $r$
and put $du/dr$ from it into~(4.2.5).  As a result we obtain
%
\begin{equation}
\frac{dp}{dr}\;\frac{(u^2 - a^2)}{p}
  = -\frac{GM}{r^2} + \frac{2u^2}{r}\,,
\tag{4.2.6}
\end{equation}
%
which shows that at the sonic point, where $a = u$,
%
\begin{equation}
2u_s^2 = \frac{GM}{r_s}\,,\qquad \text{or} \quad
u_s^2 = a_s^2 = \tfrac{1}{2}\,\frac{GM}{r_s}\,.
\tag{4.2.7}
\end{equation}

By combining this result with the Bernoulli equation~(4.2.3) we obtain the
speed of sound at the sonic point in terms of the speed of sound at infinity
%
\begin{equation}
a_s = u_s = a_\infty \left(\frac{2}{5-3\Gamma}\right)^{1/2}\!;
\tag{4.2.8}
\end{equation}
%
and using equation~(4.2.7) we obtain the radius at the sonic point
%
\begin{equation}
r_s = \left(\frac{5-3\Gamma}{4}\right)\frac{GM}{a_\infty^2}\,.
\tag{4.2.9}
\end{equation}
%
(Notice that the gravitational potential $GM/r_s$ at the sonic point is of
order~$a_\infty^2$.)

Now it is straightforward to calculate the accretion rate from
equation~(4.2.4):
%
\begin{equation}
\dot{M}_0 = 4\pi\,\alpha\,G^2 M^2\,\rho_\infty\,a_\infty^{-3}
= 4\Gamma^{3/2}\,\alpha\,G^2\,M^2\,\rho_\infty\,
  (m_p / k T_\infty)^{3/2}\,.
\tag{4.2.10}
\end{equation}

Here $\alpha$ is a constant of order unity which depends on~$\Gamma$:
%
\begin{equation}
\alpha \equiv \frac{\pi}{4\Gamma^{3/2}}
  \left(\frac{2}{5-3\Gamma}\right)^{(5-3\Gamma)/[2(\Gamma-1)]}\!.
\tag{4.2.11}
\end{equation}
%
\begin{align}
\alpha &\simeq 1.5  \quad \text{for} \quad \Gamma = 1,    \notag \\
\alpha &\simeq 1.2  \quad \text{for} \quad \Gamma = 1.4,  \notag \\
\alpha &\simeq 0.3  \quad \text{for} \quad \Gamma = 5/3;
\end{align}
%
and we have used the equation of state
$p_\infty = (\rho_\infty / m_p)\,k\,T_\infty$.

This expression for the accretion rate is the main result of our calculation.

Notice that $\dot{M}_0$ depends on~$\Gamma$.  At a typical temperature of
$T_\infty \sim 10^4$\,K the interstellar gas will be partially ionized, so
when it is compressed some energy will go into further ionization.  This means
that the value of $\Gamma$ outside and near the sonic point will be somewhat
below $5/3$.  A value of $\Gamma = 1.4$ might be reasonable for use in
equations~(4.2.10) and~(4.2.11).

It is now straightforward to derive from equations~(4.2.3) and~(4.2.4) the
radial distributions of all quantities.  The general picture is as follows.
Outside the ``radius of influence''
%
\begin{equation}
r_i = 2\,GM/a_\infty^2 = r_g\,(c/a_\infty)^2
\tag{4.2.12}
\end{equation}
%
\[
\simeq 10\,r_g \quad \text{for} \quad \Gamma = 1.4,
\]
%
the pull of the hole hardly makes itself felt, so $\rho$ and $a$ are nearly
equal to their values at ``infinity''; $\rho \simeq \rho_\infty$,
$a \simeq a_\infty$; and the density and sound velocity begin to rise.
After the gas passes through the sonic point~$r_s$, it is in near free-fall
%
\begin{equation}
u = (2GM/r)^{1/2} = a_\infty\,(r_i/r)^{1/2}
\qquad \text{at} \quad r < r_s\,;
\tag{4.2.13a}
\end{equation}
%
so the law of mass conservation requires the density to increase as
%
\begin{equation}
\rho = \frac{\dot{M}_0}{4\pi r^2\,u}
\simeq 0.2\,\rho_\infty\,(r_i/r)^{3/2}
\qquad \text{at} \quad r < r_s\,.
\tag{4.2.13b}
\end{equation}

For adiabatic compression, temperature rises as
$T \propto \rho^{\,\Gamma-1}$; consequently, at $r < r_s$
%
\begin{equation}
T \simeq 0.5\,T_\infty\left(\frac{r_i}{r}\right)^{\frac{3}{2}(\Gamma-1)}
\qquad \text{if} \quad \Gamma = 1.4~\text{near sonic point.}
\tag{4.2.13c}
\end{equation}

The above solution, (4.2.10)--(4.2.13), for adiabatic hydrodynamic accretion
is due originally to \citet{Bondi1952}.

Rewrite the accretion rate~(4.2.10) for hydrodynamic flow in a form similar
to that (eq.~4.1.3) for noninteracting particles
%
\begin{equation}
\dot{M}_0 = \pi\,r_g^2\,\rho_\infty\,c\,(c/a_\infty)^3.
\tag{4.2.14}
\end{equation}

Direct comparison, and use of the approximate equality between speed of sound
$a_\infty$ and proton speeds~$v_\infty$, shows that the hydrodynamic
accretion rate is larger by a factor $(c/a_\infty)^2 \simeq 10^{10}$ than
the accretion rate for independent particles.  The physical reason for this
is clear: gas is distinguished from independent particles by the frequent
collisions of its atoms; these collisions limit the tangential velocities
during infall, but permit radial velocities to grow.

Other, useful forms for the hydrodynamic accretion rate are
%
\begin{equation}
\dot{M}_0 \simeq 10^{5}
\left(\frac{M}{M_\odot}\right)^{\!2}
\left(\frac{\rho_\infty}{10^{-24}\;\text{g\,cm}^{-3}}\right)
\left(\frac{a_\infty}{10\;\text{km\,sec}^{-1}}\right)^{\!-3}
\;\text{,}
\tag{4.2.15a}
\end{equation}
%
\begin{equation}
\dot{M}_0 \simeq \left(1\times 10^{11}\;\frac{\text{g}}{\text{sec}}\right)
\left(\frac{M}{M_\odot}\right)^{\!2}
\left(\frac{\rho_\infty}{10^{-24}\;\text{g\,cm}^{-3}}\right)
\left(\frac{T_\infty}{10^4\;\text{K}}\right)^{\!-3/2}\!.
\tag{4.2.15b}
\end{equation}

%%% ====================================================================
%%% 4.3  Thermal Bremsstrahlung from the Accreting Gas
%%% ====================================================================

\subsection{Thermal Bremsstrahlung from the Accreting Gas}
\label{sec:4.3}

The above idealized case of spherical, adiabatic, hydrodynamic accretion can be
complicated by various physical processes.  Consider, first, the effects of energy
loss due to thermal bremsstrahlung.  To calculate the bremsstrahlung most easily,
we shall assume that the energy loss has a negligible effect on the thermal energy
density of the gas, and then we shall check that this assumption is valid.

When the effects of bremsstrahlung losses are neglected, then the temperature,
density, and velocity distributions retain their adiabatic forms (4.2.13).  One can
readily check that for $T_\infty \sim 10^4$\,K and $\Gamma = 1.4$ the gas remains nonrelativistic,
$kT \ll m_p c^2$ (but just barely so) down to $r \simeq r_s$.  Consequently, the rate at which
one gram of gas at radius $r$ radiates thermal bremsstrahlung is [cf.\ eq.~(2.1.29)]
%
\begin{equation}
\varepsilon_{ff}
= \bigl(5 \times 10^{20}~\text{ergs/(g\,sec)}\bigr)\,
  (\rho / \text{g\,cm}^{-3})
  \left(\frac{T_\infty}{10^4\,\text{K}}\right)^{\!1/2}
  \left(\frac{r_i}{r}\right)^{\!3(\Gamma-1)/2}
\tag{4.3.1}
\end{equation}
%
For comparison, the rate at which adiabatic compression increases the energy of
one gram of gas is [eq.~(2.6.41)]
%
\begin{equation}
\varepsilon_{\text{ad.\ heating}}
= \frac{p}{\rho}\,\frac{d\rho}{\rho\,dt}
= -\frac{3}{2}\,\frac{u}{r}\,\frac{p_\infty}{\rho_\infty}
  \left(\frac{r_i}{r}\right)^{3\Gamma/2}\,.
\tag{4.3.2}
\end{equation}
%
Using expression (4.2.12) for $r_i$, and the thermodynamic relations for a hydro\-gen gas
%
\begin{equation}
a_\infty^2 = (\Gamma p_\infty / \rho_\infty) = \Gamma kT_\infty / m_p\,,
\tag{4.3.3}
\end{equation}
%
we can bring this heating rate into the form
%
\begin{equation}
\varepsilon_{\text{ad.\ heating}}
= \bigl(3 \times 10^4~\text{ergs}/(\text{g\,sec})\bigr)
  \left(\frac{T_\infty}{10^4\,\text{K}}\right)^{\!5/2}
  \left(\frac{M}{M_\odot}\right)^{\!-1}
  \left(\frac{r_i}{r}\right)^{3\Gamma/2}\,.
\tag{4.3.4}
\end{equation}

This heating rate greatly exceeds the free-free loss rate (4.3.1) at all radii.  Hence,
we are justified in our use of the adiabatic forms of $T(r)$, $\rho(r)$, and $u(r)$.  The total
power radiated as thermal bremsstrahlung is
%
\begin{equation}
L_{ff}
= \int_{r_g}^{\infty} \varepsilon_{ff}\,\rho\, 4\pi r^2\, dr
\sim \left(5 \times 10^{31}~\frac{\text{ergs}}{\text{sec}}\right)
  \left(\frac{M}{M_\odot}\right)^{\!3}
  \left(\frac{T_\infty}{10^4\,\text{K}}\right)^{\!-3.3}
  \left(\frac{\rho_\infty}{10^{-24}~\text{g\,cm}^{-3}}\right)^{\!2}\,,
  \quad\text{for }\Gamma = 1.4.
\tag{4.3.5}
\end{equation}

Perhaps a more realistic value for $\Gamma$ would be a little less than $5/3$ down to
$r \sim 10^3 r_g$ at which point $kT \sim m_e c^2$, and then $\Gamma = 4/3$ below that radius.  In
this case $L_{ff}$ would be increased by several orders of magnitude [cf.\ the rela%
tivistic corrections in eq.~(2.1.31)].  Even with such an increase, and even for
$M = 100\,M_\odot$, $L_{ff}$ is so small that it is of little or no observational interest.  For
this reason, we shall not attempt a more rigorous calculation of it.

%%% ====================================================================
%%% 4.4  Influence of Magnetic Fields and Synchrotron Radiation
%%% ====================================================================

\subsection{Influence of Magnetic Fields and Synchrotron Radiation}
\label{sec:4.4}

Up to now we have ignored the fact that the interstellar gas possesses a magnetic
field.  For interstellar temperatures of interest, $T_\infty \sim 10^4$\,K, the magnetic field
will be frozen into the gas (cf.\ \S\ref{sec:2.5}).  The strength of the typical intergalactic
field, $B_\infty \sim 10^{-6}$\,G, is such that its energy density and pressure are not far below
those of the gas
%
\begin{equation}
\frac{B_\infty^2}{8\pi}
\sim 4 \times 10^{-14}~\frac{\text{ergs}}{\text{cm}^3}\,.
\tag{4.4.1}
\end{equation}

At large radii the field may be small enough that one can ignore its influence on
the accretion.  In this case the standard hydrodynamical inflow will stretch each
fluid element out radially (cross-sectional area of fluid element $\propto r^2$, hence
diameter $\propto r$; volume $\propto \rho^{-1} \propto r^{3/2}$, hence radial diameter $\propto r^{-1/2}$).
This radial stretch and tangential compression must quickly convert the initial
field in the fluid element into a nearly radial field with strength
%
\begin{equation}
B \propto \text{(cross sectional area)}^{-1} \propto r^{-2}
\tag{4.4.2}
\end{equation}
%
(conservation of flux).  Hence, the magnetic energy in a given fluid element will
rise as
%
\begin{equation}
E_{\mathrm{mag}}
= \frac{B^2}{8\pi}~\text{(volume of element)}
\propto r^{-4}\,\rho^{-1}
\propto r^{-4}\,r^{3/2}
= r^{-5/2}\,.
\tag{4.4.3}
\end{equation}

Not long after the gas crosses the sonic point, this magnetic energy will become
so large as to exceed the thermal energy in the fluid element
%
\begin{equation}
E_{\mathrm{therm}}
= \frac{3\,kT}{2\,m_p}~\text{(mass of element)}
\propto r^{-(3/2)(\Gamma-1)}\,.
\tag{4.4.2a}
\end{equation}

At this point magnetic pressures must come into play, and the flow will cease
to obey the standard adiabatic hydrodynamic laws of \S\ref{sec:4.2}.

The form of the modified magnetohydrodynamic flow is not understood at
all well today.  The best guesses and calculations to date are those of Shvartsman
(1971).  They suggest a turbulent flow, with reconnection of field lines between
adjacent ``cells'' (cf.\ \S\ref{sec:2.9}), and with a rough equipartition of the gravitational
potential energy between magnetic fields, kinetic infall of gas, turbulent energy,
and thermal kinetic energy of gas particles:
%
\begin{equation}
\bigl\{\text{gravitational energy}\bigr\}
= \frac{GM}{r}
\sim 4u^2
\sim 4\rho v_{\mathrm{turb}}^2
\sim 4\,\frac{3kT}{m_p}
\sim \frac{B^2}{8\pi\rho}
\tag{4.4.2b}
\end{equation}
%
per unit mass.

The mass density corresponding to this equipartition condition,
$4\pi r^2 \rho u = \dot{M}_0$, because $\dot{M}_0$ is determined by conditions
outside and at the sonic point, where the magnetic field is not yet large enough
to strongly influence the flow, we can use equation~(4.2.15b) for $\dot{M}_0$ and write
%
\begin{equation}
\rho
\sim \left(6 \times 10^{-12}~\frac{\text{g}}{\text{cm}^3}\right)
  \left(\frac{\rho_\infty}{10^{-24}~\text{g\,cm}^{-3}}\right)
  \left(\frac{T_\infty}{10^4\,\text{K}}\right)^{\!-3/2}
  \left(\frac{r_i}{r}\right)^{\!3/2}\,.
\tag{4.4.3a}
\end{equation}

Accepting this educated guess as to the nature of the flow, we can ask what
types of radiation might be emitted.  As before~(\S\ref{sec:4.3}) bremsstrahlung is too weak
to be interesting.  However, synchrotron radiation can be rather strong.  Most of
the synchrotron radiation will come from the high-temperature, strong-field
region near the Schwarzschild radius.  There $kT \gg m_e c^2$, so the relativistic
formula~(2.3.13) for synchrotron radiation is applicable,
%
\begin{equation}
\varepsilon_{\mathrm{synch}}
\simeq \bigl(2.2 \times 10^{-10}~\text{ergs}/(\text{g\,sec})\bigr)\,
  \bar{r}_e^2\, B_\perp^2\,;
\tag{4.4.4}
\end{equation}
%
and, because the gas turns out to be optically thin, the total luminosity is
%
\begin{equation}
L_{\mathrm{synch}}
\simeq \int_{r_g}^{\infty}\varepsilon_{\mathrm{synch}}\,\rho\, 4\pi r^2\, dr
\sim \bigl(1 \times 10^{32}~\text{ergs}/\text{sec}\bigr)
  \left(\frac{M}{M_\odot}\right)^{\!2}
  \left(\frac{\rho_\infty}{10^{-24}~\text{g\,cm}^{-3}}\right)^{\!2}
  \left(\frac{T_\infty}{10^4\,\text{K}}\right)^{\!-3/2}\,.
\tag{4.4.5}
\end{equation}

For comparison, the rest mass-energy being accreted is
%
\begin{equation}
\dot{M}_0 c^2
\sim \bigl(1 \times 10^{32}~\text{ergs}/\text{sec}\bigr)
  \left(\frac{M}{M_\odot}\right)^{\!2}
  \left(\frac{\rho_\infty}{10^{-24}~\text{g\,cm}^{-3}}\right)
  \left(\frac{T_\infty}{10^4\,\text{K}}\right)^{\!-3/2}\,.
\tag{4.4.6}
\end{equation}

Since most of the radiation comes from $r \sim 2r_g$ where the thermal energy is
${\sim}\,3$ to 10 per cent of the rest mass-energy, the synchrotron losses account for
${\sim}\,3$ to 10 per cent of the thermal energy available.  However,
for $M \gtrsim 10\,M_\odot$ the effects of the energy losses on the temperature and thence
on the luminosity will not exceed a factor of ${\sim}\,1$.

Equations~(4.4.5) and (4.4.6) predict that a hole of $M \sim 10\,M_\odot$ will convert
${\sim}\,1$ per cent of the rest mass of the infalling gas into synchrotron radiation,
producing a luminosity of ${\sim}\,10^{32}$\,ergs/sec, which is about 3 per cent of the
luminosity of the sun.  The radiation, located at [cf.\ eq.~(2.3.18)] $r \simeq 2r_g$, will have a spectrum
with a broad maximum located at
%
\begin{equation}
\nu_{\mathrm{peak}}
\sim (7 \times 10^{14}~\text{Hz})
  \left(\frac{\rho_\infty}{10^{-24}~\text{g\,cm}^{-3}}\right)^{\!1/2}
  \left(\frac{T_\infty}{10^4\,\text{K}}\right)^{\!-3/4}
\tag{4.4.7}
\end{equation}
%
($\nu \sim 7 \times 10^{14}$\,Hz, $\lambda \sim 4000$\,\AA); and it will die out exponentially for
$\nu \gtrsim 7 \times 10^{14}$\,Hz; cf.\ \S\ref{sec:2.3}).  As Shvartsman (1971) remarks, such a spectrum
and luminosity resemble those of ``DC white-dwarf stars'' (white dwarfs with
continuous, line-free spectra).  This has led Shvartsman to speculate that
``perhaps some of the objects heretofore regarded as type~DC white dwarfs are
actually black holes.''

Instabilities and inhomogeneities in the flow (due, e.g., to reconnection of
field lines) might lead to marked fluctuations in the luminosity of the hole.  The
timescales for such fluctuations might be of the order of the gas travel time from
${\sim}\,10\,r_g$ to ${\sim}\,2r_g$; i.e.
%
\begin{equation}
\Delta t_{\mathrm{fluctuations}}
\sim (10^{-3}~\text{to}~10^{-4}~\text{sec})(M/M_\odot)\,.
\tag{4.4.8}
\end{equation}

Such rapid fluctuations in an object so faint should be very difficult to detect,
even with the 200-inch telescope.  For further details on the fluctuations, see
the end of \S\ref{sec:4.7}.

%%% ====================================================================
%%% 4.5  Interaction of Outflowing Radiation with the Gas
%%% ====================================================================

\subsection{Interaction of Outflowing Radiation with the Gas}
\label{sec:4.5}

As the luminosity $L$ pours out from the region $r \sim 2r_g$, it interacts with the
inflowing gas.  For the luminosities $L \sim 10^{32}$\,ergs/sec and frequencies $\nu \sim 10^{14}$\,Hz
derived in the last section, one can readily verify that the infalling gas is optically
thin to free-free absorption and to synchrotron reabsorption.  (See \S\ref{sec:2.6} for
basic concepts and equations used in such a calculation.)

What of electron scattering?  Electron scattering can exert 3 types of influence:
(i)~If the electron concentration is sufficiently high, it can act as a ``blanket'' to
impede the outflow of the radiation.  To test for such blanketing one can examine
the optical depth of the gas to electron scattering.
%
\begin{equation}
\tau_{es}(r = 2r_g)
= \int_{2r_g}^{\infty}\kappa_{es}\,\rho\, dr
\simeq (1 \times 10^{-9})(M/M_\odot)
  (\rho_\infty / 10^{-24}~\text{g\,cm}^{-3})
  (T_\infty / 10^4\,\text{K})^{-3/2}\,.
\tag{4.5.1}
\end{equation}

(A more careful calculation would replace the nonrelativistic absorption
coefficient $\kappa_{es} = 0.4$\,cm$^2$/g by the coefficient appropriate to relativistic
electrons, since $T \sim 10^{12}$\,K at $r \sim r_g$; such a calculation would also reveal
extreme optical thinness.)  (ii)~Since the outgoing photons have energies
$h\nu \sim 10$\,eV far less than the electron kinetic energies, Compton scattering can
modify the spectrum, producing high-energy photons [see eq.~(2.4.9)].  However,
the extreme optical thinness guarantees that only about one photon in a million
scatters; so this effect can be ignored.  (iii)~The outpouring photons exert a
pressure on the infalling gas when they scatter, and this effect can retard the
inflow.  Notice that the magnitude of the first two effects is governed by the
density of the inflowing gas.  By contrast, the magnitude of the pressure effect
is governed by the luminosity of the outpouring radiation.  Let us calculate this
effect explicitly, at radii sufficiently large that the electrons are nonrelativistic.

Because the electron scattering cross-section $d\sigma/d\Omega = \tfrac{1}{2}r_0^2(1+\cos^2\theta)$ is
symmetric between forward and backward angles, each scattered photon transfers
its full momentum, $h\nu/c$, to the electron.  Hence, the time-averaged
photon force acting on an electron at radius $r$ is
%
\begin{equation}
\langle\text{Force}\rangle = \left\langle\frac{d(\text{momentum})}{dt}\right\rangle
= \int \frac{d(\text{number of photons})}{d(\text{area})\,dt\,d\nu}\, h\nu\,\sigma_e\, d\nu
\tag{4.5.2}
\end{equation}
%
\[
= \sigma_e F = \sigma_e (L/4\pi r^2)\,,
\]
%
where $\sigma_e = (8\pi/3)r_0^2 = 0.657 \times 10^{-24}$\,cm$^2$ is the total scattering cross section,
$F$ is the radiation flux (ergs\,cm$^{-2}$\,sec$^{-1}$) and $L$ is the total luminosity of the hole.
Notice that this time-averaged force is completely independent of the spectrum
of the radiation (so long as most of the photons have $h\nu \ll m_e c^2$ as seen in the
electron rest frame, so that nonrelativistic cross sections are applicable).  It is
also independent of the optical thickness---being equally valid for $\tau \gg 1$, in
which case $F \ll J$ (see \S2.6 for notation).

The pull of gravity must work against this $\sigma_e (L/4\pi r^2)$ force.  The photon force
acts almost entirely on the electrons ($\sigma_e \sim 1$\,(mass)$^{-2}$, so each proton feels a
force $3 \times 10^6$ times weaker than that felt by an electron).  By contrast, the
acceleration of gravity acts equally on electrons and protons.  Hence, the electrons
move outward slightly relative to the protons, creating an electric field that
transmits the photon force $\sigma_e L/4\pi r^2$ from electrons to protons.  The net force
that then acts on each proton (assuming a completely ionized hydrogen gas) is
%
\begin{equation}
\langle\text{Force}\rangle_{\text{total}} = \frac{\sigma_e L}{4\pi r^2} - \frac{GMm_p}{r^2}\,.
\tag{4.5.3}
\end{equation}

Thus, there is a critical luminosity [``Eddington'' (1926) limit]
%
\begin{equation}
L_{\text{crit}} \equiv \frac{4\pi G M m_p}{\sigma_e}
= (1.3 \times 10^{38}~\text{ergs/sec})\,(M/M_\odot)
\tag{4.5.4}
\end{equation}

such that for $L \ll L_{\text{crit}}$ gravity dominates at all radii, but for $L \gg L_{\text{crit}}$ photon
pressure dominates at all radii.  Eddington (1926) derived this limit for the case of
stellar equilibrium.  Zel'dovich and Novikov (1964) derived this limit for the
cases of quasars and accretion.

For the case of current interest---a black hole of $M\sim 1$ to $100\,M_\odot$ swallowing
interstellar gas---the luminosity is far below the critical value, so the photon
pressure can be ignored.

However, in other contexts the accretion rate $\dot{M}_0$ may be so great, and the
efficiency $\xi$ for converting the inflowing mass-energy into outgoing luminosity
$L$ may be sufficiently high that
%
\begin{equation}
\frac{L}{L_{\text{crit}}} = \frac{\xi \dot{M}_0 c^2}{L_{\text{crit}}}
\lesssim \frac{\xi}{0.1} \left(\frac{\dot{M}_0}{10^{18}~\text{g/sec}}\right)
\tag{4.5.5}
\end{equation}

may temporarily exceed unity.  When this happens, the photon pressure will
impede further mass inflow, thereby reducing the luminosity to $L \lesssim L_{\text{crit}}$.  Such
accretion is said to be ``\emph{self-regulated}''.  Self-regulated accretion, with modifications due to lack of spherical symmetry, may be important for black holes in
close binary systems; see \S5.13.  However, it is unimportant for the case at
hand.

The outflowing radiation from a black hole in interstellar space ($L \sim 10^{32}$
ergs/sec, $\nu \sim 10^{14}$--$10^{15}$\,Hz) can also interact with the surrounding interstellar
gas.  One can verify that the radiation is sufficiently intense to keep the gas at a
temperature $T_\infty \sim 10^4$\,K in the neighborhood of $r = r_i$, and to keep it partially
ionized.

%%% ====================================================================
%%% 4.6  Validity of the Hydrodynamical Approximation
%%% ====================================================================

\subsection{Validity of the Hydrodynamical Approximation}
\label{sec:4.6}

The above model for accretion onto an interstellar black hole relies crucially on
the validity of the hydrodynamical approximation at and outside the sonic point,
and on equipartition assumptions inside the sonic point.  If the gas were to behave
like independent particles rather than like a fluid, then the accretion rate would
be greatly reduced (see \S4.1).  Thus, it is crucial to check the validity of the
hydrodynamical approximation.

In an ionized hydrogen plasma \emph{without} magnetic fields, the deflection of
protons away from straight-line paths is caused primarily by the cumulative
effects of many small-angle Coulomb scatterings.  The total distance a proton
must travel before the ``random-walk'' deflections add up to an angle of ${\sim}90^\circ$
is [see, e.g., Shkarofsky et al.\ (1966) or Pikel'ner (1961)]
%
\begin{equation}
\lambda_p = \frac{9}{4\pi}\,\frac{(kT)^2}{n_p e^4 L_c}
\simeq (7 \times 10^{12}~\text{cm})\left(\frac{\rho_\infty}{10^{-24}~\text{g\,cm}^{-3}}\right)^{-1}
\left(\frac{T}{10^4\,\text{K}}\right)^{2}\,.
\tag{4.6.1}
\end{equation}

Here $L_c$ is a ``Coulomb-logarithm'' (analogue of Gaunt factor in theory of
bremsstrahlung) given by
%
\begin{equation}
L_c \simeq 23 + \tfrac{3}{2}\ln (n_e/\text{cm}^{-3});
\tag{4.6.2}
\end{equation}

$n_e$ and $n_p$ are the number densities of electrons and (ionized) protons; and we
have evaluated $\lambda_p$ for the case of a fully ionized hydrogen plasma.  The plasma
will behave like a fluid on scales $l \gg \lambda_p$, but like independent particles on scales
$l \ll \lambda_p$.

The scale of interest for determining the accretion rate is the sonic radius
%
\begin{equation}
r_s = \frac{5 - 3\Gamma}{4}\,\frac{GM}{a_\infty^2}
\simeq (10^{13}~\text{cm})\left(\frac{M}{M_\odot}\right)
\left(\frac{T_\infty}{10^4\,\text{K}}\right)^{-1}\,,
\tag{4.6.3}
\end{equation}

at which point $T \sim T_\infty$, $\rho \sim \rho_\infty$, so
%
\begin{equation}
\lambda_p(r_s)/r_s \sim 0.7 \left(\frac{M}{M_\odot}\right)^{-1}
\left(\frac{T_\infty}{10^4\,\text{K}}\right)^{-1}\,.
\tag{4.6.4}
\end{equation}

Thus, if one ignores the magnetic field, then the gas will not quite behave like a
fluid near the sonic point---and it may fail even more to be fluid-like at small
radii; cf.\ \S\,13.4 of ZN.

The interstellar magnetic field, of strength $B_\infty \sim 10^{-6}$\,G, can ``save'' the
hydrodynamic approximation.  It deflects protons away from straight-line
motion in a distance of the order of the Larmour radius
%
\begin{equation}
\lambda_L = \frac{m_p c\,v}{eB}
\simeq (1 \times 10^{8}~\text{cm})
\left(\frac{B}{10^{-6}~\text{G}}\right)^{-1}
\left(\frac{T}{10^4\,\text{K}}\right)^{1/2}
\tag{4.6.5}
\end{equation}

and this is true even when the thermal pressure of the plasma exceeds the
magnetic pressure so that the field gets dragged about by the gas.  The Larmour
radius is small compared to all macroscopic scales of interest (e.g.\ small compared
to $r_s$ when one uses the values $B \sim B_\infty$ and $T \sim T_\infty$ appropriate to $r_s$).  Hence,
the hydrodynamical approximation is fairly well justified.

%%% ====================================================================
%%% 4.7  Gas Flow Near the Horizon
%%% ====================================================================

\subsection{Gas Flow Near the Horizon}
\label{sec:4.7}

Near the horizon, $r \sim r_g$, relativistic effects will modify the Newtonian picture
of accretion.  It is in fact difficult to solve for the relativistic flow using the
relativistic Bernoulli equation (2.5.33) and the law of mass conservation (2.5.23$'$).
Such a solution is analogous to the Newtonian hydrodynamical solution studied
in \S4.2.  However, such a solution is partially irrelevant because it ignores the
role of the magnetic field, and such a solution is not needed because the intensity
of the gravitational pull, as viewed in stationary frames near the horizon, is so
great as to guarantee supersonic flow with qualitatively the same form as free-particle
fall.  Thus, from a knowledge of geodesic orbits one can infer the main
features of the flow.

Consider, first, spherical accretion onto a nonrotating hole (Schwarzschild
gravitational field).  A fluid element in the accreting gas, like a freely falling
particle, reaches the horizon in finite proper time.  The moment of crossing the
horizon is in no way peculiar, as seen by proper time.  The density does not
reach infinity there; and the temperature does not reach infinity.  In fact,
because the equation for radial geodesics in the Schwarzschild geometry
%
\begin{equation}
ds^2 = (1 - r_g/r)\,c^2\,dt^2 + (1 - r_g/r)^{-1}\,dr^2 + r^2\,d\Omega^2
\tag{4.7.1}
\end{equation}

has the first integral
%
\begin{equation}
(dr/d\tau)^2 = 2GM/r \quad\text{(for fall from rest at } r \gg r_g\text{)},
\tag{4.7.2}
\end{equation}

and because the law of rest-mass conservation $\boldsymbol{\nabla}\cdot(\rho \mathbf{u}) = 0$ has the first integral
%
\begin{equation}
4\pi r^2 \sqrt{-g}\, \rho\,(dr/d\tau) = \dot{M}_0\,,
\tag{4.7.3}
\end{equation}

the density of rest mass must vary as
%
\begin{equation}
\rho = \frac{\dot{M}_0}{4\pi r^2} \left(\frac{r}{2GM}\right)^{1/2}
\simeq \left(6 \times 10^{-12}~\text{g\,cm}^{-3}\right)
\left(\frac{M}{M_\odot}\right)^{-1}
\left(\frac{T_\infty}{10^4\,\text{K}}\right)^{-3/2}
\,(r/r_g)^{-3/2}\,.
\tag{4.7.4}
\end{equation}

This must be as true near and inside the horizon, $r \lesssim r_g$, as it is in the Newtonian
realm $r \gg r_g$.  Similarly, the temperature will retain its Newtonian form (4.4.2b):
%
\begin{equation}
T \simeq (10^{12}~\text{K})(r_g/r)\,;
\tag{4.7.5}
\end{equation}

and the synchrotron emissivity will retain its Newtonian value (4.4.4).  However,
the radiation which reaches infinity will be sharply cut off as the gas element
passes through the horizon.
%
\begin{equation}
\left(\text{fraction of emitted radiation that}\atop
\text{reaches } r = \infty \text{ from infalling gas at radius } r_f\right)
\simeq \frac{27}{64}\left(1 - \frac{r_g}{r_f}\right)\,.
\tag{4.7.6}
\end{equation}

In previous sections we have taken rough account of this cutoff and of the
accompanying redshift by retaining the Newtonian luminosity down to $r = 2r_g$.

In the case of a rotating (Kerr) hole, near and inside the ergosphere the
dragging of inertial frames will swing the infalling gas into orbital rotation
about the hole.  As the gas approaches the horizon, its angular velocity as seen
from infinity must approach the angular velocity of the horizon,
%
\begin{equation}
\Omega \to \Omega_{\text{horizon}} = \frac{a}{r_+^2 + a^2}
= \frac{c^3}{2GM}
\quad\text{for ``maximally rotating hole'', } a = M
\tag{4.7.7}
\end{equation}
%
(geometrized units, $c = G = 1$).

Here $a$ is the hole's angular momentum per unit mass.  Any extra-luminous hot
spot in the gas (e.g.\ a region where magnetic field lines are reconnecting; cf.\
\S2.9) must orbit the hole with this angular velocity as it falls inward.  The
brightness of the hole as seen at Earth will be modulated with a period
$P \simeq 4\pi r_g/\Omega \gtrsim (10^{-4}~\text{sec})(M/M_\odot)$ as a result of the Doppler shift of the spot's
light and the bending of its rays.  (A factor of $2\pi/\Omega$ is the orbital period; the
other factor of $2\pi/\Omega$ is the light travel time between the radius at which the
spot is located at the beginning of one period and the radius to which it has
fallen at the end of the period.)  Thus, one might expect quasiperiodic
fluctuations in the light from a black hole living alone in the interstellar medium.

%%% ====================================================================
%%% 4.8  Accretion onto a Moving Hole
%%% ====================================================================

\subsection{Accretion onto a Moving Hole}
\label{sec:4.8}

Turn attention now from a hole at rest in the interstellar medium, to a hole that
moves with a speed much larger than the sound speed, $u_\infty \gg a_\infty$. Examine the
resulting accretion in the rest frame of the hole. At radii $r \gg 2GM/u_\infty^2$, the kinetic
energy per unit mass, $\frac{1}{2}u_\infty^2$, is far larger than the potential energy, $GM/r$; so the
pull of the hole has no effect on the gas flow. At $r \sim 2GM/u_\infty^2$, the pull of the
hole becomes significant. Because the flow is supersonic the gas will respond to
that pull in the same manner as would noninteracting particles; its pressure
cannot have an influence. (One can see this, for example, in the Euler equation
%
\begin{equation}
\text{const.} = \tfrac{1}{2}v^2 + \frac{a^2}{\Gamma - 1} - \frac{GM}{r}\,,
\tag{4.8.1}
\end{equation}

where the speed-of-sound (enthalpy) term---which accounts for pressure effects---is
negligible compared to the kinetic-energy term.) Thus, the flow lines will be
identical to the trajectories of test particles: they will be hyperbolae
(Fig.~\ref{fig:4.8.1}a).

After passing around the hole, the trajectories of test particles intersect
(particle collisions; dotted line of Figure~\ref{fig:4.8.1}a). In the gas-dynamic case such
intersection of flow lines is prevented by rapidly mounting pressure. Consequently,
a shock front must develop around the black hole (Fig.~\ref{fig:4.8.1}b).
Outside the shock front the flow lines will be hyperbolic test-particle trajectories.
Behind the shock front, they will not. The shock will be located roughly where
potential energy equals kinetic energy, so its characteristic size $l_s$ as shown in
Figure~\ref{fig:4.8.1}b will be
%
\begin{equation}
l_s \sim GM/u_\infty^2\,.
\tag{4.8.2}
\end{equation}

In the shock, the gas will lose most of its velocity perpendicular to the shock
front [``strong shock'' in terminology of \S\,2.7; ``strong'' because $u_\infty/a_\infty ={}$
(speed of gas)/(speed of sound) $\gg 1$ on front side]; but it will retain all of its
velocity parallel to the front (denote this by $u_\|$). For a given gas element, if the
remaining kinetic energy behind the front greatly exceeds the potential energy,
$\frac{1}{2}u_\|^2 \gg GM/r$, then the gas element will escape the pull of the hole. If
$\frac{1}{2}u_\|^2 \lesssim GM/r$, then it cannot escape. The dividing line between escape and
capture is a dotted line in Figure~\ref{fig:4.8.1}b; and the corresponding impact parameter
is labelled $b_\text{capture}$.

We can calculate $b_\text{capture}$ in order of magnitude by idealizing the shock as
confined to the thin dotted region of Figure~\ref{fig:4.8.1}a (Hoyle and Lyttleton~1939).
The value of $b_\text{capture}$ will correspond to an orbit for which
%
\begin{equation}
\tfrac{1}{2}v_\perp^2 = GM/x \quad \text{at}\quad y = 0.
\tag{4.8.3}
\end{equation}

Straightforward examination of Kepler orbits reveals that
%
\begin{equation}
b_\text{capture} = 2GM/u_\infty^2\,.
\tag{4.8.3$'$}
\end{equation}

Since the mass flux in the gas at large radii is $\rho_\infty u_\infty$, the accretion rate is
%
\begin{equation}
\dot{M}_0 = (\pi b_\text{capture}^2)\,\rho_\infty\, u_\infty
= 4\pi G^2 M^2 \rho_\infty / u_\infty^3\,.
\tag{4.8.4}
\end{equation}

\begin{figure}[htbp]
\centering
% \includegraphics[width=\columnwidth]{fig_4_8_1}
\fbox{\parbox{0.9\columnwidth}{\centering [Figure 4.8.1 placeholder]\\[6pt]
(a)~Trajectories of test particles (and of supersonic gas)
relative to a black hole, as viewed in the rest frame of the hole.\\[4pt]
(b)~The flow lines of supersonic gas
relative to a black hole, as viewed in the rest frame of the hole.
The trajectories and flow lines are virtually the same, until
the gas encounters the shock front.}}
\caption{\textit{Figure~4.8.1.} The trajectories of test particles~(a), and the flow lines
of supersonic gas~(b) relative to a black hole, as viewed in the rest frame of the hole.
The trajectories and flow lines are virtually the same, until the gas encounters the
shock front.}
\label{fig:4.8.1}
\end{figure}

More careful calculations give accretion rates which differ from this by
multiplicative factors of ${\sim}\,0.5$ to $1.0$ (Bondi and Hoyle~1944).

Notice that the accretion rate~(4.8.4) for the supersonic case, and the rate
(4.2.10) for a hole at rest are identical except for the replacement of sound
speed $a_\infty$ by speed of flow $u_\infty$, and except for a numerical coefficient of order
unity. In the intermediate case of $u_\infty \sim a_\infty$, one can use the hybrid formula
%
\begin{equation}
\dot{M}_0 \simeq \frac{4\pi G^2 M^2 \rho_\infty}{(a_\infty^2 + u_\infty^2)^{3/2}}
\tag{4.8.5}
\end{equation}
%
\begin{equation}
\dot{M}_0 \simeq \left(1 \times 10^{11}~\frac{\text{g}}{\text{sec}}\right)
\frac{(M/M_\odot)^2\;(\rho_\infty/10^{-24}~\text{g~cm}^{-3})}
{\bigl[(u_\infty/10~\text{km~sec}^{-1})^2 + (a_\infty/10~\text{km~sec}^{-1})^2\bigr]^{3/2}}\,.
\notag
\end{equation}
%
(Bondi~1952, Hunt~1971).

[Salpeter~(1964) gave the first application of these formulas to black holes; all
previous work dealt with normal stars.]

It is worth noting that a significant fraction of the kinetic energy released in
the shock
%
\begin{equation}
\frac{dE}{dt} \simeq \tfrac{1}{2}\dot{M}_0\, u_\infty^2
\simeq \left(10^{24}~\text{ergs~sec}^{-1}\right)
\left(\frac{M}{M_\odot}\right)^{\!2}
\left(\frac{\rho_\infty}{10^{-24}~\text{g~cm}^{-3}}\right)
\left(\frac{u_\infty}{10~\text{km~sec}^{-1}}\right)^{-1/2}
\tag{4.8.6}
\end{equation}
%
may go into excitation of unionized atoms and thence into light as the atoms
decay, thus producing a shock front that glows (albeit weakly). See, e.g.,
Pikel'ner~(1961), and \S\S\,9 and~10 of Kaplan~(1966).

The motion of the gas behind the shock is practically unknown, particularly
when one takes account of magnetic fields. Many different types of instabilities
may occur, both here and in the case of a stationary black hole. For example at
radii $r < l_s$, where the inflow has become supersonic again, turbulence may
create shock waves which convert kinetic energy into heat (Bisnovatyi-Kogan and
Sunyaev~1971). A variety of plasma instabilities may also arise. And the
reconnection of magnetic field lines will surely be important. One can only guess that
the overall, time-averaged flow will resemble the equipartition model of
Shvartsman~(\S\,4.4).

One of the factors that may be important near the gravitational radius, where
most of the luminosity is produced, is angular momentum. Suppose that the accreting
gas has nonzero angular momentum per unit mass, $\mathscr{L}$, in the direction of motion
of the hole. If the angular momentum per unit mass, $\mathscr{L}$, exceeds
$r_g c$, then centrifugal forces will become important before the infalling gas reaches
the horizon; the gas will be thrown into circulating orbits; and only after viscous
stresses have transported away the excess angular momentum will the gas fall
down the hole. (See \S\,5 for a situation similar to this.) The resulting viscous
heating and lengthened time spent outside the hole, as well as the changed flow
pattern, may affect the outpouring luminosity significantly. Moreover, when the
hole moves from a region with angular momentum of gas in one direction to a
region with angular momentum in another, the gas flow may become particularly
violent for a while near the horizon; and a strong flare of luminosity may be
produced (Salpeter~1964; Shvartsman~1971). Almost nothing is known today
about these issues.

Is the specific angular momentum of the accreting gas likely to exceed $r_g c$?
Interstellar gas is accreted from a region of diameter
%
\begin{equation}
2b_\text{capture} = \frac{2r_g}{u_\infty^2/c^2}
= (3 \times 10^{14}~\text{cm})
\left(\frac{M}{M_\odot}\right)
\left(\frac{u_\infty}{10~\text{km~sec}^{-1}}\right)^{-2}\,.
\tag{4.8.7}
\end{equation}

(For $u_\infty < a_\infty$, we must replace $u_\infty$ by $a_\infty$.) The interstellar gas is turbulent with
the turbulent velocity $v_\text{turb}$ on scales $l$ given roughly by
%
\begin{equation}
v_\text{turb} = \left(10^5~\frac{\text{cm}}{\text{sec}}\right)
\left(\frac{l}{3 \times 10^{20}~\text{cm}}\right)^{q}\,,
\tag{4.8.8}
\end{equation}
%
where $q$ is a coefficient believed to lie between $1/2$ and~$1$ (Kaplan and Pikel'ner
1970), and the coefficient $10^6$\,cm/sec is taken from astronomical observations.
Let us take $q = 0.75$ as a best guess. Then the specific angular momentum
associated with this turbulence, for a gas element of size $l$, is
%
\begin{equation}
\mathscr{L}_\text{turb} \simeq v_\text{turb}\, l
\simeq \left(3 \times 10^{26}~\frac{\text{cm}^2}{\text{sec}}\right)
\left(\frac{l}{3 \times 10^{20}~\text{cm}}\right)^{1.75}\,.
\tag{4.8.9}
\end{equation}

Consequently, the specific angular momentum of the gas that accretes from a
region of size $l = 2b_\text{capture}$ should be
%
\begin{equation}
\frac{\mathscr{L}_\text{turb}}{r_g c}
\simeq 1 \times
\left(\frac{M}{M_\odot}\right)^{3/4}
\left(\frac{u_\infty}{10~\text{km~sec}^{-1}}\right)^{-7/2}\,.
\tag{4.8.10}
\end{equation}

Evidently, the angular momentum may be important in some cases and may not
be in others.

%%% ====================================================================
%%% 4.9  Optical Appearance of Hole: Summary
%%% ====================================================================

\subsection{Optical Appearance of Hole: Summary}
\label{sec:4.9}

In summary, a black hole alone in the interstellar medium may be expected to
emit ${\sim}\,10^{29}$ to $10^{35}$\,ergs/sec of synchrotron radiation, concentrated in the
optical part of the spectrum (${\sim}\,10^{14}$\,Hz to $10^{15}$\,Hz). The output may
fluctuate on timescales of ${\sim}\,10^{-2}$ to $10^{-4}$ seconds due to instabilities, and to
orbital motion of hot spots about the hole (\S\S\,4.4, 4.7, 4.8). There may also
be strong flares when the hole moves out of one turbulent cell of interstellar
gas into another. The time interval between flares would be about
%
\begin{equation}
\Delta t_\text{flares}
\simeq 2b_\text{capture}/u_\infty
\simeq (10~\text{years})
\left(\frac{M}{M_\odot}\right)
\left(\frac{u_\infty}{10~\text{km~sec}^{-1}}\right)^{-3}\,.
\tag{4.9.1}
\end{equation}

It will probably be difficult observationally to distinguish accreting, isolated
black holes from accreting, old neutron stars without magnetic fields. The only
distinguishing feature will be the emission from the surface of the neutron star.
